%===============================================================================
% Zweck: KTR-Seminar-Vorlage in Anlehung an G. Wirtz, Lehrstuhl Praktische Informatik
%===============================================================================
%===============================================================================
% zentrale Layout-Angaben und Befehle
%===============================================================================
%
% \usepackage{german}
\usepackage{ifthen}
% Language Selection for Babel
\ifthenelse{\equal{\lang}{ngerman}}{\usepackage[german, \lang]{babel}}{\usepackage[\lang]{babel}}

\usepackage[utf8]{inputenc}
\usepackage{fancyhdr}
\usepackage[T1]{fontenc}
\usepackage{ae}
\usepackage{color}
\usepackage{amsmath}
\usepackage{amsfonts}
%

\usepackage{caption}

%%   Fuer anspruchsvolle Tabellen   %%
\usepackage{longtable, colortbl}
\usepackage{multicol, multirow}

\usepackage[pdftitle={\@title},pdfauthor={\@author},pdftex,bookmarksopen,bookmarksnumbered,pdfborder=0]{hyperref}
\usepackage[pdftex]{graphicx}
\usepackage{float}
\usepackage{tikz}
\usepackage{pgfplots}
\usetikzlibrary{arrows,shapes,fit,positioning,snakes,backgrounds,shadows}

\pdfcompresslevel=9

% Code-Hervorhebung
% Quellcode
\usepackage{verbatim}            % Quellcode einbinden (\verbatiminput) standardpaket
\usepackage{moreverb} 
% PseudoCode
\usepackage{algorithm}
\usepackage{algpseudocode}
%\usepackage{algorithmicx}
\floatname{algorithm}{Algorithmus}
\algrenewcommand{\algorithmiccomment}[1]{\hskip1em\textcolor{gray!60}{$\rhd$ #1}}
\ifthenelse{\equal{\lang}{ngerman}}{%
\renewcommand{\listalgorithmname}{Algorithmen}
\def\algorithmautorefname{Algorithmus}
}{%
\renewcommand{\listalgorithmname}{List of Algorithms}
\def\algorithmautorefname{Algorithm}
}


%%   intoc zur Aufnhame des Abkuerzungs- und Symbolverzeichnisses ins Inhaltsverzeichnis  
\usepackage[intoc]{nomencl}
\setlength{\nomlabelwidth}{.20\hsize}
%\renewcommand{\nomlabel}[1]{#1 \dotfill}
\setlength{\nomitemsep}{-\parsep}
\makenomenclature
\ifthenelse{\equal{\lang}{ngerman}}{%
\renewcommand{\nomname}{Abk\"urzungsverzeichnis}
%\newcommand{\nomaltname}{Symbolverzeichnis}
}{%
\renewcommand{\nomname}{List of Abbreviations}
%\newcommand{\nomaltname}{List of Symbols}
}%
%\newcommand{\nomaltpreamble}{}
%\newcommand{\nomaltpostamble}{}
%\newcommand{\usetwonomenclatures}{\nomenclature[\switchnomitem]{}{}}
%\newcommand{\switchnomitem}{R}
%\renewcommand{\nomgroup}[1]{%
%\ifthenelse{\equal{#1}{\switchnomitem}}{\switchnomalt}{}}
%\newcommand{\switchnomalt}{%
%\end{thenomenclature}
%\newpage
%\renewcommand{\nomname}{\nomaltname}
%\renewcommand{\nompreamble}{\nomaltpreamble}
%\renewcommand{\nompostamble}{\nomaltpostamble}
%\begin{thenomenclature}
%}


%%   Hervorhebung der Abkuerzungsbuchstaben   %%
\usepackage[normalem]{ulem}
\newcommand{\m}[1]{\uline{#1}}

%\usepackage[dvips]{rotating}
%
% ausf\"{u}hrlichere Fehlermeldungen
\errorcontextlines=999
%
% Page-Layout: A4 aus Header
% Alternative
\setlength\headheight{14pt}
\setlength\topmargin{-15,4mm}
\setlength\oddsidemargin{-0,4mm}
\setlength\evensidemargin{-0,4mm}
\setlength\textwidth{160mm}
\setlength\textheight{252mm}
%
%% Absatzeinstellungen
\setlength\parindent{0mm}
\setlength\parskip{2ex}


\makeatletter
\renewcommand{\maketitle} {
  \begin{titlepage}
  \centering
    \begin{minipage}[t]{16cm}
      \hfill
      \begin{minipage}{12cm}
      \ifthenelse{\equal{\lang}{ngerman}}{%
        \centering
        Otto-Friedrich-Universit\"at Bamberg
        \\[12pt]%
        {\Large Professur f\"ur Informatik\\[.5em]%
        \large insbesondere Kommunikationsdienste,\\[.5em]%
        Telekommunikationsdienste und Rechnernetze}%
        }{%
        \centering
        Otto-Friedrich-University of Bamberg
        \\[12pt]
        {\Large Professorship for Computer Science,\\[.5em]%
        \large Communication Services, Telecommunication\\[.5em]%
        Systems and Computer Networks}%
        }
      \end{minipage}
      \hfill
      \begin{minipage}{3cm}
        \includegraphics[height=28mm]{config/images/logo} %height=26mm
      \end{minipage}
    \end{minipage}\\[70pt]%[50pt]
   % 
    {\Large\bf \ifthenelse{\equal{\lang}{ngerman}}{Ausarbeitung des KTR-Projekts}{Project on}}
    \\[36pt]
   % oder: {\LARGE Ausarbeitung des KTR-Seminars}\\[12pt]
    {\LARGE \@title}\\[80pt]
%    {\Large\bf \ifthenelse{\equal{\lang}{ngerman}}{Thema:}{Topic:}}\\[36pt]
%    {\LARGE\bf \subtitle}\\
    \vfill
    \begin{minipage}{\textwidth}
      \center
      \ifthenelse{\equal{\lang}{ngerman}}{Vorgelegt von:}{Submitted by:}\\
      {\Large \@author \\[18pt]}
      \ifthenelse{\equal{\lang}{ngerman}}{Betreuer:}{Supervisor:} \supervisor \\[12pt]
      Bamberg, \@date\\
      \semester
    \end{minipage}
  \end{titlepage}
}
\makeatother
%
\usepackage[intoc]{nomencl}
\ifthenelse{\equal{\lang}{ngerman}}{%
\renewcommand{\nomname}{Abkürzungsverzeichnis}}%
{\renewcommand{\nomname}{List of Abbreviations}}
\renewcommand{\nomlabel}[1]{#1 \dotfill}
\makenomenclature

% Einbindung eines Bildes mit angegebener Breite
% #1 = label f\"{u}r \ref-Verweise
% #2 = Name des Bildes ohne Endung relativ zu Bilder-Verzeichnis
% #3 = Beschriftung
% #4 = Breite des Bildes im Dokument in cm
\newcommand{\bildw}[4]{%
  \begin{figure}[htb]%
    \centering
    \includegraphics[width=#4cm]{Bilder/#2}%
    \vskip -0.3cm%
    \caption{#3}%
    \vskip -0,2cm%
    \label{#1}%
  \end{figure}%
}
%
% Einbindung eines Bildes mit Seitenbreite
% #1 = label f\"{u}r \ref-Verweise
% #2 = Name des Bildes ohne Endung relativ zu Bilder-Verzeichnis
% #3 = Beschriftung
\newcommand{\bild}[3]{%
  \begin{figure}[htb]%
    \centering%
    \includegraphics[width=\textwidth]{Bilder/#2}%
    \vskip -0.3cm%
    \caption{#3}%
    \vskip -0,2cm%
    \label{#1}%
  \end{figure}%
}
%
\numberwithin{equation}{section}
%
%===============================================================================
% zentrale Layout-Angaben und Befehle
%===============================================================================
%
%#1 Breite
%#2 Datei (liegt im image Verzeichnis)
%#3 Beschriftung
%#4 Label fuer Referenzierung
\newcommand{\image}[4]{%
\begin{figure}[H]%
\centering%
\includegraphics[width=#1]{image/#2}%
\caption{#3}%
\label{#4}%
\end{figure}%
}

%#1 Datei (liegt im graphic Verzeichnis)
%#2 Beschriftung
%#3 Label fuer Referenzierung
%#4 Skalierungsfaktor
\newcommand{\scaletikzimage}[4]{%
\begin{figure}[H]%
\centering%
\scalebox{#4}{%
\input{graphic/#1.tikz}}%
\caption{#2}%
\label{#3}%
\end{figure}
}

%#1 algorithm name
%#2 algorithm label
%#3 file name in code-folder
\newcommand{\pseudo}[3]{%
\small%
\begin{algorithm}[H]%
\caption{#1}%
\label{#2}%
\input{code/#3.tex}%
\end{algorithm}%
\normalsize%
}

